\honest{Joy in Discovery}{Eliezer Yudkowsky}{2008}

I open with Lagrange: Newton was the greatest genius, and the most fortunate, ``for we cannot find more than once a system of the world to establish.''

I enjoy discovering things for myself more than reading about them. Discovering what no one else knows is something more. There is a story of a physicist who told his girlfriend that he was ``the only man in the world who knows why they shine.'' Numerous sources, I say, attest that being first to solve a great mystery is a tremendous high. \nb{The post names none.}

What bothers me is the exclusivity. Why should a discovery be worth less because someone else knows the answer? I offer three readings. The charitable one: months of struggle give an understanding that falls into place all at once. The less charitable one: we value what cost us more, as in the studies of fraternity pledges and of wine in pricier bottles. \nb{The classic initiation study, by Aronson and Mills (1959), used women joining a discussion group, not fraternity pledges. In the wine study the tasters were only told a price, which was not a cost they bore, so it does not illustrate this reading.} The uncharitable one: status. I strongly suspect, too, that people think free knowledge cannot be important.

If the joy must be scarce, only one person per civilization can have it for each truth, and Newton used up much of Earth's physics fun. Worse, ``in the Standard Model of physics'' the universe is spatially infinite, inflationary and branching, so aliens or ``Tegmark duplicates of Newton'' may have found gravity first. \nb{The Standard Model is the theory of particle physics and says none of this. The three claims are Tegmark's first three levels of parallel universes; the argument needs only infinite space, which Tegmark counts as part of the standard cosmological model.} Someone somewhere knows every answer, so the requirement to be first leads to angst that cannot be resolved, ``and I regard that as a reductio.'' \nb{This reaches only the view that no one anywhere may know. The opening story, and the post's later remark on rewards to one's civilization, are about being first on one Earth.}

My solution is to treat mystery as personal. If you do not know how your hand closes into a fist, you are as ignorant as a hunter-gatherer, whoever else knows. Telling your civilization something new is a separate reward, one-shot per discovery, like a Nobel Prize.

So I restore magic and mystery to everything you do not personally understand. Think how much you do not know. I end with Ran Prieur, on a friend's remark that Prieur looked at the world as if seeing it for the first time: ``Wait! I never have seen it before!''
